\subsection{Secants, compression, and the reference cost}
\label{paper:ep-reference}
Throughout this appendix the ambient alphabet has at most $H$ letters;
all constants may depend on $H$.  For $z\geq1$ put
\[
 E_q(z)=(q+z)\log(q+z)-z\log z,\qquad E_0(z)=0.
\]
The following elementary estimate supplies both the qualitative
compression and the quantitative separation needed below.
\begin{lemma}[Normalized secants]\label{paper:ep-secants}
For $1\leq b<y\leq H$, let
\[
 L(b,y)=\frac{y\log y-b\log b}{y-b},\qquad
 \beta_b(y)=\frac{\log L(b,y)}{\log y}.
\]
Then $\beta_b'(y)\leq-1/(24H^2)$.
\end{lemma}
\begin{proof}
Set $F(t)=\log\int_0^1e^{tu}\,du$, $0\leq t\leq\log H$.
Its second derivative is the variance under the density
$p_t(u)=e^{tu}/\int_0^1e^{tv}\,dv\geq H^{-1}$, so
\[
 F''(t)\geq H^{-1}\int_0^1(u-\mathbb E_tu)^2\,du\geq(12H)^{-1}.
\]
For $t=\log y$, $\beta_1(y)=1-F(t)/t$, whence
\[
 -\beta_1'(y)=\frac{\int_0^t vF''(v)\,dv}{yt^2}
 \geq\frac1{24Hy}\geq\frac1{24H^2}.
\]
For general $b$, writing $s=b/y$ gives
\[
 y\partial_yL(b,y)=h(s)=\frac{1-s+s\log s}{(1-s)^2},\qquad
 h'(s)=\frac{(1+s)\log s+2(1-s)}{(1-s)^3}<0.
\]
Indeed $\log s-2(s-1)/(s+1)$ increases to zero on $(0,1)$,
its derivative being $(s-1)^2/[s(s+1)^2]$.  Increasing $b$ at fixed
$y$ increases $L$ and decreases $y\partial_yL$.  It therefore decreases
\[
 \beta_b'(y)=\frac{(y\partial_yL/L)\log y-\log L}
                   {y(\log y)^2}.
\]
Comparison with $b=1$ proves the assertion.
\end{proof}

\begin{lemma}[Shifted compression]\label{paper:ep-compression}
Let $0\leq e\leq1$ and $F_r(q)=q-(r+z)^{-e}E_q(z)$.
At finitely many increasing integer capacities $r$, let integers $d_r$
be nondecreasing and satisfy $0\leq d_r\leq r$.  There is an integer
$R\geq0$ such that
\[
 F_r(\min(r,R))\leq F_r(d_r)\quad\hbox{at every capacity}.
\]
The assertion remains true if only the coefficient at the largest
capacity is decreased to an arbitrary positive value.  No restriction
on the increments $d_{r'}-d_r$ is required.
\end{lemma}
\begin{proof}
Read capacities in increasing order, starting with benchmark $R=0$.
At capacity $r$, replace $R=A$ by the current count $C=d_r$ only if
$F_r(A)>F_r(C)$.  Monotonicity of the counts ensures $A<C\leq r$.
The failing comparison means
\[
 L(z+A,z+C)>(z+r)^e.
\]
For an earlier capacity $A<h<C$, Lemma~\ref{paper:ep-secants} gives
$L(z+A,z+h)>(z+h)^e$; for $C\leq h\leq r$ the displayed
inequality suffices, and for $h\leq A$ the clipped arguments coincide.
Thus every earlier comparison is preserved.  If the last coefficient
is $a\leq(r+z)^{-e}$, its failing comparison instead gives
$L(z+A,z+C)>1/a\geq(z+r)^e$, which proves the same preservation.
For the singleton convention
$\phi_m(n)=n-1-m^{-e}n\log n$, $1\leq n\leq m$, the identical
argument uses $L(A,C)>m^e$; zero counts may first be replaced by one,
since the entropy margin is zero at both zero and one.
\end{proof}

Fix a maximal equal-product platform $[a,b]$ of product $\sigma$ in
the complete natural chain.  Write
\[
 r_i=b-i+1,\quad S_i=r_i+z,\quad c_i=S_i^{\sigma-1},\quad
 V_i=\sum_{l\leq i}L_l,\quad T=V_b.
\]
The permanent child has weight $z\geq1$.  Retain every original old
physical block, with its natural coefficient
\begin{equation}\label{paper:ep-old-coeff}
 \kappa_i=\frac{Z_i}{S_a}\prod_{i\leq l<a}\beta_l\quad(i<a),
 \qquad \kappa_i=c_i\quad(a\leq i\leq b).
\end{equation}
Its mean letter intensities are $\kappa_i(1,\ldots,1,z)$.
All old blocks have capacity $r_a$; native block $i$ has capacity $r_i$.
Put $K_{\rm in}=\sum_{i<a}L_i\kappa_i$.  The platform stationarity
identities give
\begin{align}
 K_{\rm in}E_{r_a}(z)-r_aV_{a-1}
    +\sum_{i=a}^bL_i[c_iE_{r_i}(z)-r_i]&=0,\label{paper:ep-balance}\\
 \Lambda_R:=RV_{a-1}-K_{\rm in}E_R(z)
    +\sum_{i=a}^bL_i\{\min(r_i,R)-c_iE_{\min(r_i,R)}(z)\}
    &=\lambda_{b-R}\geq0.\label{paper:ep-suffix-buffer}
\end{align}
Here $0\leq R\leq r_a$ and $\lambda_{a-1}=\lambda_b=0$.
For example, at $R=r_j$ the left side is precisely the incoming
buffer minus the remaining native drift, hence the multiplier
$\lambda_{j-1}$.  This also checks the two zero endpoint values.
These coefficients describe the unconditioned natural reference;
conditional estimates retain each old word's actual list.

\begin{proposition}[Reference nonnegativity]\label{paper:ep-reference-positive}
Every lawful original coarse history on this platform alphabet has
nonnegative mean total cost.  More precisely, its cost has the form
\begin{equation}\label{paper:ep-reference-certificate}
 \overline C=D_{\rm partition}+D_{\rm clock}
  +\sum_v\omega_v\left(\lambda_{b-R_v}
                  +\sum_i L'_i\Delta_{i,v}\right),
\end{equation}
where all terms on the right are nonnegative, $\sum_v\omega_v=1$,
and the finite list of vertices has cardinality bounded in terms of
$H$ alone.
\end{proposition}
\begin{proof}
First retain the chain of cells containing the permanent child,
$A_0\supsetneq\cdots\supsetneq A_q$, with split times
$\tau_0\leq\cdots\leq\tau_q$.  Resolve every pure sibling immediately
when it leaves this chain.  This is lawful, since every name in a live
cell is eligible until its split.  If $d(t)$ is the number of unresolved
boundaries, then $d$ is nonincreasing, $d(t)\leq r_i$ in native block
$i$, and the resulting star-spine cost is exactly
\begin{align}
 \overline C_{\rm sp}
 &=\overline F_{\rm full}(\tau_0)
       +\sum_{l=1}^q\{\overline F_{A_l}(\tau_l)
                         -\overline F_{A_l}(\tau_{l-1})\}\\
 &=\int_0^T[d(t)-\kappa(t)E_{d(t)}(z)]\,dt.\label{paper:ep-spine-integral}
\end{align}
The initial value is zero by \eqref{paper:ep-balance}; no pure-cell
lifetime or original rank delay remains in this comparison history.

The transport identity is
$\kappa_i=\kappa_{i+1}\beta_i^{1-s_i}$, including the old/native
transition, so $\kappa$ increases on $[0,V_a]$.  Define
\[
 \bar c_a=\frac{K_{\rm in}+L_ac_a}{V_a}\leq c_a,\qquad
 A(x)=\int_0^x(\bar c_a-\kappa(v))\,dv.
\]
Then $A\geq0$, $A(0)=A(V_a)=0$.  Since $E_{d(t)}$ is nonincreasing,
Stieltjes summation gives the exact correction
\[
 D_{\rm clock}=\int_0^{V_a}(\bar c_a-\kappa(t))E_{d(t)}\,dt
              =\sum_{0<x<V_a}\Delta E_{d}(x)A(x)\geq0.
\]
Replace this first region by $L'_a=V_a$, $c'_a=\bar c_a$; for $i>a$
keep $L'_i=L_i$, $c'_i=c_i$.  The suffix sums
\eqref{paper:ep-suffix-buffer} are unchanged.

Subdivide the ordered split-time domain at the reduced physical
endpoints $0,V_a,\ldots,V_b$.  In each closed chamber the averaged
cost is affine.  At a vertex every split time is an endpoint: otherwise
a maximal tied group of nonendpoint coordinates can be moved in both
directions without violating its order, block or deadline constraints.
Thus the cost is a convex combination, with weights $\omega_v$, of
\[
 \sum_iL'_iF_i(d_i(v)),\qquad F_i(q)=q-c'_iE_q(z).
\]
At every vertex the ranks $d_i(v)$ decrease in physical order.
Lemma~\ref{paper:ep-compression}, including its reduced largest
coefficient, supplies $R_v$ for which
$\Delta_{i,v}=F_i(d_i(v))-F_i(\min(r_i,R_v))\geq0$.
This proves \eqref{paper:ep-reference-certificate} for a star spine.
In particular the total negative proper-cell increments are at most
the root mean plus the positive proper-cell increments.  The root
mean is nonnegative: it starts and ends at zero and its slopes have
one positive-to-negative sign change, as established below.

For a general coarse history let $d(t)$ count all unresolved original
boundaries.  Its eligible unresolved rank is at most $d(t)$, the
remaining rank being original delay.  A partition of masses
$(1,\ldots,1,z)$ with unresolved rank at most $d$ has entropy at most
$E_d(z)$.  Indeed merge two nontrivial cells and split off their smallest
atom.  Convexity of $x\log x$ preserves rank and does not decrease
entropy.  Iteration leaves one nontrivial cell, and increasing any
chosen mass increases its entropy; it therefore uses the largest
$d+1$ atoms, namely the child and $d$ exterior names.  Consequently
\[
 D_{\rm partition}=\int_0^T\left\{\kappa_iE_{d(t)}(z)
             -\sum_{A\in\Pi(t)}\operatorname{Ent}_{v_i}(A)\right\}\,dt\geq0.
\]
Apply the preceding chamber argument to the history's actual activation
order to obtain \eqref{paper:ep-reference-certificate}.
\end{proof}
For a proper platform, whose $z$ is bounded away from one, the finite
partition inequality has a strict gap unless its unresolved partition
is one child-containing cell of full unresolved rank and singletons.
It has no gap between different child-containing sets of the same
cardinality.  At $z=1$ further equal-weight ties remain.  Moreover,
Proposition~\ref{paper:ep-reference-positive} concerns this platform
alphabet: old cells crossing to a name outside it belong to their
larger original context.

\subsection{Discrete reserves and single crossings}
\label{paper:ep-reserves}
For groups of capacity $c_j$, length $L_j\geq0$ and coefficient
$0<\kappa_j\leq1$, put
\[
 g_j(q)=q-\kappa_jE_q(z),\qquad
 \Lambda_q=\sum_jL_jg_j(\min(c_j,q)).
\]
Assume $c_j+z\leq H$ and all $\Lambda_q\geq0$.
\begin{lemma}[Discrete curvature reserve]\label{paper:ep-curvature}
Let $T_a=\sum_{c_j=a}L_j$ and
$W_a=\sum_{c_j>a}\kappa_jL_j$.  Then
\begin{equation}\label{paper:ep-curvature-local}
 \Lambda_a+\tfrac12\log H\,T_a\geq\frac{W_a}{2H}.
\end{equation}
If $A<R<B$, set
$G_R(a)=(\min(R,a)-A)(B-\max(R,a))/(B-A)$ for $A<a<B$.
Then
\begin{equation}\label{paper:ep-green}
 \Lambda_R+\log H\sum_{a=A+1}^{B-1}G_R(a)T_a
 \geq\frac{(R-A)(B-R)}{2H}
                  \sum_{c_j\geq B}\kappa_jL_j.
\end{equation}
In particular a native group of capacity greater than $a$ contributes
at least $L_i/(2H^2)$ to the right side of
\eqref{paper:ep-curvature-local}.
\end{lemma}
\begin{proof}
With $Q_a=E_{a+1}-2E_a+E_{a-1}$,
\[
 Q_a=\int_{-1}^1\frac{1-|t|}{z+a+t}\,dt\geq H^{-1},
\]
and direct subtraction gives
\[
 D_a:=2\Lambda_a-\Lambda_{a-1}-\Lambda_{a+1}
 =Q_aW_a+\sum_{c_j=a}L_j[1-\kappa_j(E_a-E_{a-1})].
\]
The last bracket is at least $-\log H$.  Since
$2\Lambda_a\geq D_a$, this proves \eqref{paper:ep-curvature-local}.
The discrete Green identity expresses $\Lambda_R$ as the linear
interpolation of $\Lambda_A,\Lambda_B$ plus
$\sum_{a=A+1}^{B-1}G_R(a)D_a$.  These endpoint values are nonnegative,
$W_a\geq\sum_{c_j\geq B}\kappa_jL_j$, and
$\sum_aG_R(a)=(R-A)(B-R)/2$, giving \eqref{paper:ep-green}.
A native coefficient is at least $H^{-1}$.
\end{proof}

\begin{lemma}[Weak edges and their order]\label{paper:ep-weak-order}
On one platform let $g_r(q)=q-(r+z)^{-e}E_q(z)$, $0\leq e\leq1$,
and put $\epsilon_*=(\log2)/(192H^2)$.
If $0\leq q<r$ and $|g_r(r)-g_r(q)|\leq\epsilon_*$, then for
$q<h<r$
\begin{equation}\label{paper:ep-intervening}
 g_h(q)-g_h(h)\geq\frac{\log2}{32H^2}.
\end{equation}
Two such full-input edges $r_1\to q_1$, $r_2\to q_2$ with
$r_1>r_2$ necessarily satisfy $q_1>q_2$.
\end{lemma}
\begin{proof}
Weakness gives
$|(r+z)^{-e}L(q+z,r+z)-1|\leq\epsilon_* /(r-q)$.
Its logarithm is at least $-2\epsilon_*$.  Integrating
Lemma~\ref{paper:ep-secants} between $h+z$ and $r+z$, and using
$\log(h+z)\geq\log2$, gives
\[
 \log[(h+z)^{-e}L(q+z,h+z)]
 \geq\frac{\log2}{24H^2}-2\epsilon_*
 \geq\frac{\log2}{32H^2}.
\]
Since $g_h(q)-g_h(h)=(h-q)[(h+z)^{-e}L(q+z,h+z)-1]$,
this proves the first assertion.
For the second write $\beta(q,r)=\beta_{q+z}(r+z)$.
Weakness implies $|\beta(q_j,r_j)-e|\leq2\epsilon_*/\log2$.
If $q_1\leq q_2$, increasing the lower secant endpoint and then using
Lemma~\ref{paper:ep-secants} gives
\[
 \beta(q_2,r_2)-\beta(q_1,r_1)
 \geq\beta(q_1,r_2)-\beta(q_1,r_1)\geq\frac1{24H^2},
\]
whereas the two weak bounds give at most $1/(48H^2)$.
\end{proof}
Thus surviving deepest output sets are strictly nested in physical
order within one platform.

We record explicitly the uniform old coefficients used in the estimates.
If every natural product lies in $[\varepsilon_H,1-\varepsilon_H]$,
then \eqref{paper:ep-old-coeff} and its transport identity give
\begin{equation}\label{paper:ep-old-uniform}
 H^{-H-1}=:\kappa_*\leq\kappa_l\leq1,
 \qquad \kappa_{l+1}-\kappa_l\geq d_*>0
\end{equation}
for consecutive old groups, where one may take
\[
 d_* =\min\left\{\frac{\varepsilon_H}{H^2},
 H^{-H-1}\big[1-(1-H^{-1})^{\varepsilon_H}\big]\right\}.
\]
The first term also covers the last old/native transition.  This follows
by inserting the natural transport ratio
$\kappa_l/\kappa_{l+1}=\beta_l^{1-s_l}$ and the bounds on its
finite alphabet.  In particular all old full-rank slopes decrease.
For later quantitative comparisons one may fix
\[
 \delta_*=\frac{\kappa_*}{4H},\quad
 \epsilon_0=\min\{\delta_*/2,d_*\delta_*/(4H)\},\quad
 C_{\rm len}=H/\delta_*,\quad
 c_{\rm prefix}=\frac{d_*}{2(1+C_{\rm len})}.
\]

\begin{lemma}[Single crossing and length allocation]\label{paper:ep-single-crossing}
Take any equal-product interval $[a,b]$, not necessarily maximal, with
common product $\sigma$, child size $C=h_{b+1}$ and
$z=\exp(G_{b+1}/\sigma)$; at the terminal boundary $C=z=1$.
Put $n_i=h_i-C$, $Z_i=(n_i+z)^\sigma$, and
\[
 e_i=(n_i+z)^{\sigma-1}E_{n_i}(z)-n_i,\quad
 I_{\rm in}=\sum_{l<a}L_lZ_l\prod_{l\leq v<a}\beta_v,
 \quad M_H=H(1+\log H).
\]
Then
\begin{equation}\label{paper:ep-interval-drift}
 \sum_{i=a}^bL_ie_i
 =n_aV_{a-1}-I_{\rm in}\frac{E_{n_a}(z)}{n_a+z}
       -\lambda_{a-1}+\lambda_b.
\end{equation}
Writing $B_0=M_HV_{a-1}+|\lambda_{a-1}|+|\lambda_b|$,
the absolute value of the left side is at most $B_0$.
The signs of $e_i$ have at most one negative-to-positive change.
Except for at most one index, $|e_i|\geq\gamma_H$, where
\[
 d_H=(24H^2)^{-1},\qquad \gamma_H=1-2^{-d_H/3}>0.
\]
Consequently
\begin{align}
 |e_i|L_i&\leq B_0+M_H\sum_{l\ne i}L_l,\label{paper:ep-length-one}\\
 \gamma_H\sum_{i\in S}L_i
 &\leq B_0+M_H\sum_{i\notin S}L_i\label{paper:ep-length-sign}
\end{align}
whenever $S$ has one sign and avoids the possible small-drift index.
If $R_i=V_{a-1}+\sum_{l\ne i}L_l+
(|\lambda_{a-1}|+|\lambda_b|)/M_H$, then
\[
 |e_i|\leq M_HR_i/L_i,\qquad
 |\sigma-\theta_{n_i}(z)|\leq\frac{HM_HR_i}{(\log2)L_i},
 \quad
 \theta_n(z)=1-\frac{\log(E_n(z)/n)}{\log(n+z)}.
\]
If the small-drift index is first or last, the other lengths satisfy
$\sum_{l\ne i}L_l\leq(B_0+|e_i|L_i)/\gamma_H$.
\end{lemma}
\begin{proof}
Vary the common product, differentiating the child weight as well.
The natural recursion gives $\partial_\sigma G_i=E_{n_i}(z)/(n_i+z)$
on the interval and its transported multiple before the interval.
Stationarity, with the two boundary multipliers retained, is exactly
\eqref{paper:ep-interval-drift}.  The transport relation
$Z_i\beta_i=Z_{i+1}\beta_i^{1-s_i}$ gives
$I_{\rm in}\leq Z_aV_{a-1}$; all displayed coefficients are bounded
by $M_H$, proving its absolute bound.
Now
\[
 \frac{e_i}{n_i}
 =\exp\{(\sigma-\theta_{n_i}(z))\log(n_i+z)\}-1.
\]
The normalized-secant estimate implies
$0<\theta_n(z)<1$ and
$\theta_{n'}(z)-\theta_n(z)\geq d_H(n'-n)$ for $n'>n$.
(The secant here is $L(z,z+n)=E_n/n$.)  The integer capacities decrease
in physical order, so the asserted sign pattern follows; at most one
threshold is within $d_H/3$ of $\sigma$.  Outside that interval the
last exponential has absolute difference from one at least
$\gamma_H$, giving the uniform separation.  Moving the remaining
terms in \eqref{paper:ep-interval-drift} proves
\eqref{paper:ep-length-one} and \eqref{paper:ep-length-sign}.
Finally the derivative of $e_i$ in $\sigma$ is at least
$(\log2)/H$, yielding the threshold-distance bound.  At an extreme
small-drift index all the others have one sign, giving the final claim.
\end{proof}
On a maximal platform the preceding lemma and the increasing old
coefficients show that the full-rank slopes $g_i(r_i)=-e_i$ change
from positive to negative at most once.  Their integral is zero by
\eqref{paper:ep-balance}.  This proves the root-mean assertion used in
Proposition~\ref{paper:ep-reference-positive}.

\subsection{Which weak transitions can remain}
\label{paper:ep-weak-selection}
All histories in this subsection retain the original names and deadlines.
Their pure siblings are resolved when they leave the star spine.
Let $e=1-\sigma\geq e_0>0$.  The following argument explains why a
small reference budget forces both the input and the output names.
\begin{proposition}[Native weak transitions]\label{paper:ep-native-weak}
There are $\epsilon_H,c_H>0$ such that a synchronous native weak edge
$A\to B$, of ranks $0\leq q<R\leq r_i$, satisfies
\[
 |g_i(R)-g_i(q)|\leq\epsilon_H
 \quad\Longrightarrow\quad
 \overline C_i^-,\overline C_i^+\geq c_H L_i,
\]
for every lawful endpoint prefix and suffix context, unless $A$ is the
full current eligible alphabet and $B$ is its true deepest $q$-name
suffix together with the child.  There is at most one exceptional
$q$ for a fixed full rank.  Across native blocks these exceptional
outputs are strictly ordered as in Lemma~\ref{paper:ep-weak-order}.
\end{proposition}
\begin{proof}
We first keep track of the full compression certificate on the original
groups.  Reading an endpoint rank profile backwards makes both ranks
and capacities nondecreasing.  Use the algorithm of
Lemma~\ref{paper:ep-compression}.  Old groups have the same largest
capacity and their coefficients are nonincreasing in this reading
order.  An old update $A\to D$ gives
$(E_D-E_A)/(D-A)>1/\kappa_{\rm current}$; every processed old
coefficient is at least the current one, so its benchmark value also
decreases.  The threshold is at least the natural largest-capacity
threshold and therefore preserves all preceding native comparisons.
The resulting identity is
\begin{equation}\label{paper:ep-original-compression}
 \overline C=\Lambda_K+\sum_jL_j\Delta_j,\qquad
 \Delta_j=g_j(d_j)-g_j(\min(c_j,K))\geq0.
\end{equation}
During this algorithm the benchmark \emph{value} at every already
processed group only decreases.  Splitting a group at a cut preserves
this argument, since repeated equal coefficients cause no difficulty.

Set $\eta=1/(4H^2)$, $\nu=e_0/H$, and initially take
$\epsilon\leq1/(16H^2)$.  Since $g_i''\leq-H^{-2}$, an integer
strictly between $q$ and $R$ has value greater than both weak endpoint
values by $\eta$, whereas an integer outside $[q,R]$ has value less
than both by $\eta$.  To see the latter, compare its secant with the
secant on $[q,R]$, whose absolute slope is at most
$\epsilon/(R-q)$; strong concavity separates the slopes by at least
$(R-k)/(2H^2)$ on the left (and symmetrically on the right).
Thus if $\min(r_i,K)\notin\{q,R\}$, either
\eqref{paper:ep-original-compression} is impossible or
$\Delta_i\geq\eta$, which pays $\eta L_i$.

Suppose now that this clipped value is $q$ or $R$, with $q\geq1$.
Every future rank is at most $q$.  The reverse benchmark must reach
$q$ strictly after block $i$: if its incoming value were $A<q$, its
value on block $i$ would be strictly below both weak endpoint values,
so block $i$ would not update it; subsequent updates only decrease
this value and cannot finish at either weak endpoint.  At the first
reverse crossing $A\to q$, occurring at capacity $h$, every already
processed future group of capacity $c\geq q$ whose actual rank is
less than $q$ gains at least
\begin{equation}\label{paper:ep-deadline-charge}
 g_c(A)-g_c(q)
 >(q-A)\left[\left(\frac{h+z}{c+z}\right)^e-1\right]
 \geq \frac{e_0}{H}=\nu.
\end{equation}
Here $h>c$; a later forward rank cannot rise back to $q$.
The last bound uses $h\geq c+1$ and
$\log(1+1/(c+z))\geq1/H$.  It follows that
\[
 \overline C\geq\min\{\eta L_i,\nu D_B\},\qquad
 D_B=\sum_{j>i:\ c_j\geq q,\ |B\cap E_j|<q}L_j.
\]
More strongly, in the accepted benchmark cases the full decomposition
retains $\Lambda_K+\nu D_{\rm rank}$, where $D_{\rm rank}$ sums all
future groups of capacity at least $q$ and actual rank below $q$.
This includes both expired names and premature refinements.

Keeping the multiplier now turns this side-length payment into a
current-length payment.  If $K=q$ and $B$ is wrong, its capacity-$q$
group has rank below $q$, so
$\overline C\geq\Lambda_q+\nu T_q$.  By
Lemma~\ref{paper:ep-curvature}, this is at least
\[
 c_{\rm curv}L_i,\qquad
 c_{\rm curv}=\frac{\min(1,2\nu/\log H)}{2H^2}.
\]
If $K=R<r_i$, the reverse algorithm first reaches $R$ at block $i$
or earlier, of capacity at least $r_i>R$.  The same comparison as
\eqref{paper:ep-deadline-charge} pays $\nu$ on the already processed
capacity-$R$ group, giving the same conclusion.  Old coefficients
are no greater than the corresponding natural one, as required for
this comparison.

It remains to consider a wrong output with $R=r_i=r$ and $K\geq r$.
Set
\[
 \epsilon^*=\min\{(16H^2)^{-1},(\log2)/(192H^2)\},\quad
 \gamma=\frac{\log2}{48H^2},\quad \mu=\min(\nu,\gamma),\quad
 M=2H(1+\log H).
\]
At every future capacity $q<h<r$, rank $q$ costs at least $\gamma$
against the full benchmark by \eqref{paper:ep-intervening}; rank
below $q$ costs $\nu$ by \eqref{paper:ep-deadline-charge}.
The capacity-$q$ group has rank below $q$ because $B$ is wrong.
Consequently, with $T'=\sum_{a=q}^{r-1}T_a$,
\begin{equation}\label{paper:ep-wrong-intermediate}
 \overline C\geq\Lambda_K+\mu T'.
\end{equation}
Set
\[
 W=\sum_{c_j\geq r}\kappa_jL_j,\qquad
 A_H=\frac{2H^2\log H}{\log2},\qquad
 \epsilon_H=\min\{\epsilon^*,1/(16HA_H)\}.
\]
With $\ell=L(q+z,r+z)$ and $e'=\log\ell/\log(r+z)$,
\[
 0<e'<1,\qquad |e'-e|\leq\frac{2\epsilon_H}{(r-q)\log2}.
\]
For an earlier/current capacity $c\geq r$, put $n=\min(c,K)\geq r$.
Since $\kappa_j\leq(c+z)^{-e}$, normalized-secant monotonicity yields
\begin{align*}
 g_j(q)-g_j(n)
 &\leq\kappa_j(n-q)[(n+z)^{e'}-(c+z)^e]\\
 &\leq A_H\epsilon_H\kappa_j.
\end{align*}
The last inequality uses $n-q\leq H$ and
$\partial_v(n+z)^v\leq H\log H$ on $[0,1]$; it is one-sided and
remains valid if $e'<e$.  Smaller capacities between $q$ and $r$
contribute at most $MT'$, and capacities at most $q$ contribute zero.
Thus
\[
 \Lambda_q\leq\Lambda_K+A_H\epsilon_HW+MT'.
\]
The Green estimate \eqref{paper:ep-green} on $[q-1,r]$, evaluated
at $q$, gives $\Lambda_q+\log H\,T'\geq(r-q)W/(2H)$.
Combining these inequalities and \eqref{paper:ep-wrong-intermediate}
shows
\[
 \frac{7W}{16H}\leq\Lambda_K+(M+\log H)T'
 \leq\max\{1,(M+\log H)/\mu\}\overline C.
\]
Since $W\geq L_i/H$, this pays
$c_{\rm wrong}L_i$, where
$c_{\rm wrong}=7\min\{1,\mu/(M+\log H)\}/(16H^2)$.
The case of a genuine deadline completion, with full rank $r$ on
the preceding blocks, is included (then $K\leq r$); no lower bound
on intervening side lengths has been imposed.

For a proper input $q<R<r_i$, the following time exchange pays the
current block even when its output is correct.  At the right endpoint
$b_i$ choose $c\in A\setminus B$ and an eligible $d\notin A$,
already cut at $\tau_d\leq l_i$.  Both names have deadlines at least
$b_i$.  Move only $c$ to $\tau_d$ to form one comparison history,
and move only $d$ to $b_i$ to form a second.  Both are lawful.  Between
$\tau_d$ and $b_i$ their ranks are $k(t)-1$ and $k(t)+1$, surrounding
the original $k(t)$, and remain between zero and eligible capacity.
Proposition~\ref{paper:ep-reference-positive} gives nonnegative
comparison costs, while exact subtraction gives
\[
 2\overline C-\overline C^- -\overline C^+
 =\int_{\tau_d}^{b_i}\kappa(t)
       [E_{k(t)+1}-2E_{k(t)}+E_{k(t)-1}]\,dt
 \geq L_i/H^2.
\]
The right endpoint costs at least $L_i/(2H^2)$; the left differs by
$L_i[g_i(R)-g_i(q)]$, so both cost at least $\eta L_i$.
This also covers $q=0$.  Taking the minimum of the three positive
constants proves the stated selection for every context and hence
for its finite minimum.  Strict concavity makes the weak matching
from a fixed full rank unique, and
Lemma~\ref{paper:ep-weak-order} gives the ordering across blocks.
\end{proof}

\begin{proposition}[Old weak transitions]\label{paper:ep-old-weak}
Every weak transition in an original old group of a successor
platform has reference endpoint costs at least $c_HL_i$.
For a full-input transition both costs are at least $c_HV_{a-1}$.
These assertions hold for every lawful star-history context, including
$z=1$ at a terminal platform.
\end{proposition}
\begin{proof}
Let the incoming platform have full exterior rank $r$, $S=r+z$, and
$V=V_{a-1}>0$.  Its averaged incoming coefficient is
$\kappa_{\rm av}=I_{\rm in}/(SV)$.  At the preceding strict node,
write $p\geq1$, $u<1$, $x=S^u$,
$A=E_p(x)/(p+x)$ and $\beta=x/(p+x)$.
Stationarity, with $m=-\lambda_{\rm preceding}\leq0$, gives
\[
 I_{\rm in}=\beta(pV+m)/A\leq \frac{x}{D(p,x)}V,
 \qquad D(p,x)=E_p(x)/p.
\]
The entropy integral and concavity give
\[
 D(p,S^u)\geq S^{u-1}(1+\log S)+\frac{p}{2H},
 \qquad D(p,x)\leq1+\log H.
\]
Indeed integrate $1+\log(x+t)$ for $0\leq t\leq p$, using
$\log(1+t/x)\geq t/(x+p)$, and then use the tangent bound for
$S^{1-u}$.  It follows that
\[
 \kappa_{\rm av}(1+\log S)\leq1-\delta_0,
 \qquad \delta_0=\frac{p}{2H(1+\log H)}.
\]
Put $\beta_H=1/[4H(1+\log H)^2]$, and use a weak threshold
at most $\min\{\beta_H/4,\kappa_*/(8H)\}$.
For a full-input edge $r\to q$, weakness implies
$\kappa_i\geq(1-\epsilon)/(1+\log S)$, hence
$\kappa_i-\kappa_{\rm av}\geq\beta_H$.
Old coefficients increase, so
\[
 \sum_{l<i}L_l(\kappa_i-\kappa_l)
 \geq V(\kappa_i-\kappa_{\rm av}).
\]
Move just the names in $A\setminus B$ to time zero, leaving the
future $B$ cuts fixed.  The comparison is lawful and nonnegative.
If the original transition time is $\tau$, exact subtraction is
\[
 \overline C_\tau-\overline C_{\rm early}
 =(E_r-E_q)\sum_{l<i}L_l(\kappa_i-\kappa_l)
       +[g_i(r)-g_i(q)]\tau.
\]
Since $E_r-E_q\geq1$, $\tau\leq V$, both endpoints cost at least
$\beta_HV/2$.  If instead $q<R<r$, use the same legal two-cut
exchange as in Proposition~\ref{paper:ep-native-weak}.
Now its second difference is at least $\kappa_*L_i/H$, giving
both endpoint costs at least $\kappa_*L_i/(4H)$.

For the eventual actual-word application, an admitted endpoint row
$C_v\geq c_0\overline C_v-h$, together with the original named
regularity estimate
$E_i(t)\geq-C_H\vartheta L_i-C_HK(1+L_i^\zeta)$,
converts this to $c'_HL_i-C$.  Choose $\vartheta$ first small, then
use Young's inequality on $L_i^\zeta$ ($\zeta<1$).  Nonweak edges
use their favored endpoint.  Multiple coordinates share the same
finite deterministic cost budget.
\end{proof}

\subsection{A terminal suffix without proper weak binary edges}
\label{paper:ep-terminal-suffix}
\begin{proposition}\label{paper:ep-binary-free}
In a terminal maximal platform with positive native lengths there is
a deterministic native suffix
whose total length is at least $c_HT$, and which has no proper
full-input weak binary edge.  Its earliest largest original block has
length at least $c_HT/H$.  The suffix and the tie rule depend only on
the natural profile.
\end{proposition}
\begin{proof}
Let its natural alphabet sizes be $H\geq h_1>\cdots>h_m\geq2$,
$r_i=h_i-1$, $\kappa_i=h_i^{\sigma-1}$, and
$e_i=h_i^\sigma\log h_i-(h_i-1)=-g_i(r_i)$.
The exact terminal balance gives
\[
 D:=\sum_iL_ie_i
 =nV_{\rm old}-\frac{E_n(1)}{h_1}I_{\rm in}
 \geq b_HV_{\rm old},\qquad
 b_H=\frac{1}{2H^2(1+\log H)^2}.
\]
For completeness, when $V_{\rm old}>0$ insert
$I_{\rm in}=\beta(pV_{\rm old}+m)/A$, $m\leq0$, from the
preceding proof.  The inequality
$E_p(x)/p\geq1+\log x+p/[2(x+p)]$ gives a surplus of at least
$np/[2H^2(1+\log H)^2]$ times $V_{\rm old}$.
When there is no incoming block, $D=0$.

If $r\to q$ is proper ($1\leq q<r$) and
$|g_i(r)-g_i(q)|\leq\epsilon_w\leq1/(4H^3)$, then
$-g_i''\geq H^{-2}$ and $g_i(0)=0$ imply
\[
 g_i(q)-\frac qrg_i(r)\geq\frac{q(r-q)}{2H^2}.
\]
Writing $v=g_i(r)-g_i(q)$ shows
$g_i(r)\geq qr/(2H^2)+rv/(r-q)\geq1/(4H^2)=:c_w$.
Thus every such weak block has $e_i\leq-c_w$.
The thresholds $\alpha_h$ solving $e_h(\sigma)=0$ are strictly
increasing in $h$, by Lemma~\ref{paper:ep-single-crossing}.
Choose disjoint fixed closed neighborhoods
$|\sigma-\alpha_h|\leq\tau$ in which $|e_h|\leq c_w/2$;
this is possible since $|\partial_\sigma e_h|\leq H(\log H)^2$.
Outside these neighborhoods the nonzero drifts have a uniform
absolute lower bound $g_0>0$.

If a present threshold is near $\sigma$, begin the suffix there;
otherwise begin at the first positive drift, which exists since
$D\geq0$.  The head has negative drift at most $-g_0$, while the
suffix has at most one small drift followed by positive drifts and
therefore contains no proper weak binary block.  Its drift sum obeys
\[
 D_{\rm tail}\geq b_HV_{\rm old}+g_0L_{\rm head}\geq0.
\]
Since $|e_i|\leq H(1+\log H)$, its length is a fixed fraction of
$V_{\rm old}+L_{\rm head}+L_{\rm tail}=T$.  The largest of its at
most $H$ original blocks has length at least $c_HT/H$; choosing the
earliest one resolves ties deterministically.
\end{proof}
This suffix supplies an unread reservoir for the one uncolored root
count.  It does not impose a within-block terminal path event.
